% 环面 T² 示意图 (Torus T^2)
% 描述: 展示二维环面的不同表示方法
% 数学概念: 环面、T² = S¹ × S¹、亏格1曲面、商空间构造
\documentclass[border=10pt]{standalone}
\usepackage{tikz}
\usepackage{amsmath,amssymb}
\usetikzlibrary{arrows.meta,positioning,calc,3d,patterns,decorations.pathmorphing}

\begin{document}
\begin{tikzpicture}[
    >=Stealth,
    scale=1.0,
    axis/.style={->, line width=1pt, color=black!60},
    arrow/.style={->, line width=1.2pt, >=Stealth}
]

% ============= 左侧: 三维透视图 =============
\begin{scope}[xshift=-4cm, yshift=1cm]
    % 标题
    \node[font=\bfseries] at (0, 3.5) {三维环面};
    
    % 环面主体 (使用椭圆弧绘制)
    % 外轮廓
    \draw[line width=1.5pt, color=blue!70!black] 
        (0, 0) ellipse (2.2 and 1.8);
    
    % 内孔
    \draw[line width=1.5pt, color=blue!70!black] 
        (0, 0) ellipse (1 and 0.8);
    
    % 填充颜色 (上半部分)
    \fill[blue!15, opacity=0.6] 
        (0, 0) ellipse (2.2 and 1.8);
    \fill[white] 
        (0, 0) ellipse (1 and 0.8);
    
    % 纬线 (水平圆)
    \draw[color=red!70!black, line width=0.8pt, dashed] 
        (0, 0.6) ellipse (1.85 and 0.4);
    \draw[color=red!70!black, line width=0.8pt] 
        (0, -0.6) ellipse (1.85 and 0.4);
    \node[color=red!70!black, font=\small] at (2.3, -0.5) {$b$ (纬线)};
    
    % 经线 (垂直圆) - 前半部分
    \draw[color=green!60!black, line width=0.8pt] 
        (1.6, 0) arc (0:180:0.4 and 1.4);
    % 经线 - 后半部分(虚线)
    \draw[color=green!60!black, line width=0.8pt, dashed] 
        (1.6, 0) arc (0:-180:0.4 and 1.4);
    \node[color=green!60!black, font=\small] at (2.3, 1) {$a$ (经线)};
    
    % 基本群生成元标记
    \draw[->, color=green!60!black, line width=1.5pt] (1.6, 0.3) arc (10:160:0.35 and 0.3);
    \draw[->, color=red!70!black, line width=1.5pt] (0.5, -0.6) arc (-160:160:0.6 and 0.2);
    
    % 中心点
    \fill[black] (0, 0) circle (2pt);
    \node[below right] at (0, 0) {$x_0$};
    
    % 坐标轴
    \draw[axis] (0, 0) -- (3, 0) node[right] {$x$};
    \draw[axis] (0, 0) -- (0, 2.5) node[above] {$z$};
\end{scope}

% ============= 中间: 平面矩形商空间 =============
\begin{scope}[xshift=2cm, yshift=1cm]
    % 标题
    \node[font=\bfseries] at (0, 3.5) {商空间表示};
    
    % 矩形
    \fill[yellow!15] (-2, -1.5) rectangle (2, 1.5);
    \draw[line width=1.5pt, color=blue!60!black] (-2, -1.5) rectangle (2, 1.5);
    
    % 边标识 - 水平边 (a)
    \draw[color=green!60!black, line width=2pt, ->] (-1.8, 1.5) -- (-0.2, 1.5);
    \draw[color=green!60!black, line width=2pt, ->] (0.2, 1.5) -- (1.8, 1.5);
    \draw[color=green!60!black, line width=2pt, ->] (-1.8, -1.5) -- (-0.2, -1.5);
    \draw[color=green!60!black, line width=2pt, ->] (0.2, -1.5) -- (1.8, -1.5);
    \node[color=green!60!black, above] at (0, 1.5) {$a$};
    \node[color=green!60!black, below] at (0, -1.5) {$a$};
    
    % 边标识 - 垂直边 (b)
    \draw[color=red!70!black, line width=2pt, ->] (-2, -1.3) -- (-2, 1.3);
    \draw[color=red!70!black, line width=2pt, ->] (2, -1.3) -- (2, 1.3);
    \node[color=red!70!black, left] at (-2, 0) {$b$};
    \node[color=red!70!black, right] at (2, 0) {$b$};
    
    % 说明
    \node[anchor=north, font=\small] at (0, -2.2) {$I \times I / \sim$};
    \node[anchor=north, font=\scriptsize, text width=4cm, align=center] 
        at (0, -2.8) {将对边按照箭头方向粘合};
    
    % 基本群关系
    \node[anchor=north, draw=gray!50, fill=white, rounded corners, inner sep=5pt, font=\small] 
        at (0, -4) {$\pi_1(T^2) = \langle a, b \mid aba^{-1}b^{-1} = 1 \rangle \cong \mathbb{Z}^2$};
\end{scope}

% ============= 右侧: 细胞分解 =============
\begin{scope}[xshift=6.5cm, yshift=1cm]
    % 标题
    \node[font=\bfseries] at (0, 3.5) {CW复形结构};
    
    % 0-细胞 (顶点)
    \fill[black] (0, 0) circle (4pt);
    \node[right, font=\small] at (0.2, 0) {$e^0$ (1个)};
    
    % 1-细胞 (边) - 水平
    \draw[color=green!60!black, line width=2pt, ->] (-1.5, 0) arc (180:0:1.5);
    \node[color=green!60!black, above] at (0, 1.6) {$e^1_a$};
    
    % 1-细胞 (边) - 垂直 (稍微偏移)
    \draw[color=red!70!black, line width=2pt, ->] (-0.8, -1.2) arc (-90:90:1.2);
    \node[color=red!70!black, right] at (0.5, 0) {$e^1_b$};
    
    % 2-细胞 (面)
    \draw[color=blue!60!black, line width=1pt, dashed] (0, 0) circle (2);
    \node[color=blue!60!black, font=\small] at (-1.4, -1.4) {$e^2$};
    
    % 附加映射说明
    \node[anchor=north, font=\scriptsize, text width=4cm, align=left] 
        at (0, -3) {
        粘贴映射: \\
        $\partial e^2 \mapsto aba^{-1}b^{-1}$
    };
    
    % Euler示性数
    \node[anchor=north, draw=gray!50, fill=white, rounded corners, inner sep=5pt, font=\small] 
        at (0, -4) {$\chi(T^2) = 1 - 2 + 1 = 0$};
\end{scope}

% ============= 底部: 代数拓扑信息 =============
\node[anchor=north, font=\small, draw=blue!30, fill=blue!5, rounded corners, inner sep=10pt] 
    at (4, -5) {
    \begin{tabular}{ll}
        \textbf{代数拓扑不变量:} & \\
        基本群: & $\pi_1(T^2) \cong \mathbb{Z} \oplus \mathbb{Z}$ \\
        同调群: & $H_0 = \mathbb{Z}, \; H_1 = \mathbb{Z}^2, \; H_2 = \mathbb{Z}$ \\
        Euler示性数: & $\chi = 0$ \\
        万有覆叠空间: & $\mathbb{R}^2$ \\
        亏格: & $g = 1$
    \end{tabular}
};

\end{tikzpicture}
\end{document}
